% docs/documento_oficial/cap1_definicion_proyecto.tex
\section{Definición del Proyecto}

% ─────────────────────────────────────────────────────────────────────────────
\subsection{Descripción del problema}

En el departamento del Meta (Colombia), las condiciones ambientales agresivas ---altas temperaturas, humedad elevada y particularidades de los suelos--- inciden de manera directa en el desempeño y la durabilidad del concreto empleado en obras civiles. A pesar de ello, existe una carencia de herramientas digitales locales que apoyen, de forma didáctica y con trazabilidad, la dosificación y la selección de aditivos según el contexto regional. Esta brecha se hace evidente en los laboratorios de materiales y en los semilleros de investigación de la Universidad Cooperativa de Colombia (UCC), donde la replicabilidad de los ensayos y la sistematización de resultados son esenciales para la formación y la transferencia de conocimiento.

La ausencia de software especializado para apoyar el diseño y la verificación de mezclas adaptadas a condiciones locales puede derivar en decisiones inconsistentes, menor durabilidad y mayores costos de mantenimiento. De acuerdo con lineamientos nacionales, es prioritario impulsar la innovación tecnológica en el sector construcción para elevar la calidad y sostenibilidad de la infraestructura (Departamento Nacional de Planeación, 2018). A su vez, gremios y estudios académicos resaltan la necesidad de incorporar herramientas digitales que integren particularidades ambientales y estándares técnicos en el diseño de mezclas para optimizar recursos y asegurar la vida útil de las estructuras (CCI, 2020; González et al., 2019).

En este contexto, se identifica la necesidad de una solución \textbf{dirigida a los estudiantes del semillero de investigación y a quienes participan en las prácticas de laboratorio de materiales de la UCC}, que les permita: (i) registrar de manera estructurada las condiciones de diseño, (ii) asistir la dosificación base y la selección de aditivos con proyecciones de resistencia a distintas edades, y (iii) garantizar la trazabilidad de los ensayos y ejercicios académicos. Es importante precisar que \textbf{DIAPCE no está concebida como una herramienta de uso general ni para profesionales en ejercicio fuera del entorno universitario}; su alcance es intencionalmente formativo e institucional. DIAPCE responde a esta necesidad mediante una aplicación multiplataforma desarrollada en Flutter que opera sin conexión (SQLite), orientada exclusivamente a estudiantes del semillero y los laboratorios de la UCC para apoyar decisiones basadas en datos, estandarizar prácticas académicas y facilitar la comparación reproducible de resultados. La herramienta no sustituye el criterio del docente ni el marco normativo, sino que lo complementa con evidencia sistematizada y accesible en el contexto regional del Meta.

\subsubsection*{Pregunta de investigación}
¿Cómo por medio de la implementación de un software se pueden realizar cálculos de optimización de mezclas cementantes a partir de un modelo matemático que use propiedades físicas, mecánicas y térmicas de concreto estructural?

% ─────────────────────────────────────────────────────────────────────────────
\subsection{Objetivos}

\subsubsection{Objetivo general}
Desarrollar un software que implemente un modelo matemático para el ajuste de proporciones de materiales cementantes en mezclas de concreto estructural, orientado a los estudiantes del semillero de investigación y de los laboratorios de materiales de la UCC, sede Meta.

\subsubsection{Objetivos específicos}
\begin{enumerate}
    \item Implementar una base de datos que permita el almacenamiento de características físicas, mecánicas y/o ambientales de mezclas de concreto estructural.
    \item Desarrollar un software que implemente un modelo matemático para las mezclas de concreto estructural con una arquitectura multiplataforma y almacenamiento local.
    \item Realizar pruebas y análisis de desempeño del programa referente a criterios de calidad.
\end{enumerate}

% ─────────────────────────────────────────────────────────────────────────────
\subsection{Justificación}

\subsubsection*{Por qué surge esta propuesta}

La necesidad de DIAPCE no nace de una carencia tecnológica abstracta, sino de una realidad concreta y repetida en el laboratorio de materiales y en el semillero de investigación de la UCC, sede Meta: cada cohorte de estudiantes parte prácticamente desde cero. Los registros de ensayos anteriores no se encuentran en un formato consultable, las condiciones de cada experimento se documentan de forma dispersa y los resultados acumulados por el grupo no se traducen en un saber colectivo que oriente las decisiones de dosificación siguientes. Cuando el conocimiento generado en prácticas sucesivas no se sistematiza, la curva de aprendizaje se reinicia con cada nueva generación de estudiantes, y la ventaja de pertenecer a un semillero activo se diluye.

A esto se suma el contexto ambiental del Meta: las condiciones de temperatura y humedad del territorio no son las que suponen las tablas genéricas de dosificación derivadas de normas internacionales. Diseñar mezclas de concreto estructural en esta región exige contrastar esas referencias con evidencia experimental propia, producida y validada localmente. Sin una herramienta que centralice y haga consultable ese acervo, cada decisión de dosificación se apoya en criterios individuales difíciles de comparar y reproducir.

\subsubsection*{Qué busca el proyecto}

DIAPCE busca convertir el conocimiento experimental acumulado por el semillero en un recurso colectivo, trazable y reproducible. No se trata de reemplazar el juicio del docente ni de automatizar el diseño de mezclas; se trata de ofrecer a los estudiantes una plataforma donde los datos de sus propios ensayos ---resistencias medidas, condiciones ambientales registradas, aditivos empleados--- se conviertan en la base real de las recomendaciones de dosificación.

En términos concretos, el proyecto busca: (i) centralizar los registros de diseño y los resultados de laboratorio en un repositorio estructurado y persistente; (ii) proporcionar asistencia en la selección de proporciones y aditivos a partir de los datos experimentales del semillero, con proyecciones de resistencia a distintas edades; y (iii) establecer una práctica de registro sistemático que permita comparar resultados entre sesiones, cohortes y condiciones ambientales distintas. Todo ello dentro de un entorno de uso exclusivamente formativo, sin pretensión de reemplazar herramientas de uso profesional o industrial.

\subsubsection*{Beneficios esperados}

Para el estudiante, DIAPCE representa un punto de contacto directo entre la teoría del diseño de mezclas y la evidencia que él mismo produce en el laboratorio. La herramienta hace visible el impacto de las variables ambientales y de la selección de aditivos sobre la resistencia del concreto, reforzando competencias que van más allá de la memorización de fórmulas: análisis crítico de resultados, formulación y contraste de hipótesis, y cultura de registro.

Para el semillero como grupo, el beneficio principal es la trazabilidad: la posibilidad de recuperar, comparar y discutir resultados de ensayos anteriores como si formaran parte de una misma conversación científica continua. Esto fortalece la identidad investigativa del grupo y eleva la calidad de las discusiones académicas entre cohortes.

Para los docentes, DIAPCE ofrece un instrumento de evaluación objetiva: los registros quedan almacenados con sus parámetros de entrada, lo que permite contrastar el proceso de decisión de cada estudiante con los resultados obtenidos y con el comportamiento del conjunto de datos regional.

\subsubsection*{Componente comunitario y aporte a la comunidad académica}

El valor de DIAPCE trasciende el uso individual. Al operar sobre un repositorio compartido de datos experimentales generados por los propios estudiantes de la UCC en el contexto del Meta, la herramienta construye ---práctica a práctica--- un patrimonio de conocimiento regional. Cada ensayo registrado enriquece la base desde la que se producen las recomendaciones futuras, lo que significa que el esfuerzo de investigación de una cohorte beneficia directamente a las siguientes.

Este modelo de acumulación progresiva del conocimiento es especialmente relevante en una región donde las condiciones ambientales particulares hacen insuficientes las referencias externas y donde la producción de datos locales tiene un valor formativo e institucional difícil de sustituir. DIAPCE propone que el semillero de investigación de la UCC deje de ser un espacio donde el conocimiento se genera y se pierde por ciclos, y se convierta en un espacio donde ese conocimiento se acumula, se discute y se refina de manera continua.

% ─────────────────────────────────────────────────────────────────────────────
\subsection{Metodología}

Como metodología de trabajo se optó por usar \textbf{FDD (Feature-Driven Development)} debido a su enfoque ágil y centrado en el cliente. FDD facilita la división del trabajo en pequeñas funciones, permitiendo una detección temprana de errores y una rápida adaptación a los cambios (Lucidchart, 2024). El tipo de investigación es mixto, con un enfoque en el desarrollo de software.

\begin{figure}[H]
    \centering
    \includegraphics[width=1\textwidth]{../images/metodolgia.jpg}
    \caption{Diagrama metodológico del proyecto con la estructura de la documentación y desarrollo.}
    \label{fig:metodologia}
\end{figure}

Las fases metodológicas adoptadas son:
\begin{enumerate}
    \item \textbf{Planteamiento:} Identificación del problema, definición de objetivos y revisión de antecedentes.
    \item \textbf{Diseño:} Modelado de la base de datos, arquitectura de la aplicación y especificación de requisitos funcionales.
    \item \textbf{Desarrollo:} Implementación en Flutter/Dart con almacenamiento local SQLite e integración del conjunto de datos experimental.
    \item \textbf{Pruebas:} Validación funcional de la interfaz, pruebas de calidad del modelo de predicción y revisión de criterios normativos.
\end{enumerate}

Las \textbf{variables de investigación} consideradas son:
\begin{itemize}
    \item \textbf{Variable independiente:} El modelo matemático y los parámetros de entrada (temperatura, humedad, relación a/c, aditivo).
    \item \textbf{Variable dependiente:} Los resultados generados por el software (proyecciones de resistencia a 7/14/28 días) y las métricas de calidad (MAE, RMSE, MAPE, F1).
    \item \textbf{Variables intervinientes:} Errores en los datos de entrada, desbalance de clases en el dataset.
    \item \textbf{Variables controladas:} Formatos de entrada de datos, entorno de pruebas y criterios de desempeño definidos.
\end{itemize}
